\section{Introduction}\label{sec:intro}

For Bayesian neural networks there is a useful low-level picture of generalization, close to Occam learning over small circuits \cite{nngen}. In the strong-prior limit, and at mean-field or stronger scaling ($\gamma\ge1$), the posterior learns any target computed by a small circuit, with sample cost polynomial in the input dimension and the circuit's size. Below that scaling ($\gamma<1$) it can need exponentially many samples. This paper asks whether gradient descent admits a comparably concrete, discrete description, and how close it is to the Occam picture.

The literature already contains the qualitative ingredients. In the rich regime, from small initialization, gradient flow passes near a sequence of saddles. Each escape adds roughly one unit or one rank, so training resembles a greedy algorithm that grows a network one feature at a time. This has been made precise, to different degrees, in several settings:
\begin{itemize}
\item proved for diagonal linear networks \cite{pesme2023saddle}, and under assumptions for matrix factorization \cite{li2021greedy};
\item proved for the first escape in deep linear networks \cite{jacot2021saddle} and deep ReLU networks \cite{bantzis2025saddle}, with the later saddles conjectured;
\item proposed as a general description of two-layer networks, ``alternating gradient flows'' \cite{kunin2025agf}.
\end{itemize}
For sparse Boolean targets, the \emph{leap complexity} of Abbe, Boix-Adser\`a and Misiakiewicz \cite{abbe2023leap,abbe2022msp} predicts that online SGD on a two-layer network needs on the order of $d^{\max(L-1,1)}$ steps. Here $L$ is the largest number of new coordinates any Fourier monomial adds when the monomials are taken in the best order. This matches correlational-SQ (CSQ) lower bounds \cite{abbe2023leap}. For Gaussian single-index models with information exponent $k$, plain online SGD has the analogous rate $d^{k-1}$ \cite{benarous2021online}. Evolution by mutation and selection is equivalent to CSQ learning \cite{valiant2009evolvability,feldman2008evolvability}, and noisy neuroevolution is equivalent on average to gradient descent \cite{whitelam2021neuroevolution}. So the original idea behind this work, ``a single circuit, mutated, keeping improvements'', is already placed: it is CSQ learning.

What we add is quantitative and empirical. We train rich-regime networks on 41 Boolean targets (21 designed, 20 random formulas) at up to four input dimensions. We then build discrete models that reproduce the networks' sample complexity target by target, out of sample, and use ablations to find which ingredient of the model matters. We also prove a set of ``training has not taken long, so the circuit is still small, so a union bound over small circuits applies'' theorems, and test what can actually be certified for the trained networks.

\paragraph{Main findings.}
\begin{enumerate}
\item \textbf{Gradient descent pays a search-noise cost, not a description-length cost} (\cref{sec:networks}). The networks' $\nstar$ grows like $d^{1.04}$, $d^{1.91}$ and $d^{2.40}$ for 2-, 3- and 4-bit parities (\cref{tab:dscale}). An Occam learner needs about $k\log d$ samples.
\item \textbf{A greedy, noise-limited growth model matches the networks} (\cref{sec:model}).
  \begin{itemize}
  \item \emph{The rule.} A candidate feature is accepted once $n\rho^2$ exceeds the noise level of the continuous search that would find it. Here $\rho$ is its correlation with the residual, and the noise level grows by about a factor $d/a$ for each new input beyond the second.
  \item \emph{Simulation.} The race model, a simulation of feature growth from which the rule was distilled, matches $\nstar$ on 20 fresh random formulas within $1.15\times$ on average. The rule itself, simulated as a minimal model, is within about $1.4\times$.
  \item \emph{Closed form.} Evaluated in closed form along the target's greedy path, the rule is within about $1.34\times$ on the same formulas, and it predicted the 3-parity at $d=120$ before the networks were run.
  \item \emph{Ablations.} The race dynamics we originally built into the model matter little: on held-out targets, removing them moves $\nstar$ by at most about 0.1 beyond the model's own spread, with one exception. The finite-sample noise of the direction search does matter.
  \end{itemize}
\item \textbf{Memorize like a kernel, generalize like a feature learner} (\cref{sec:hybrid}; a \status{Heuristic} reading of \status{Numerical} results). On partially learned cells, the networks' mistakes are best predicted, among the models we tested, by the NNGP kernel of their architecture. A model that adds a kernel background to the grown features needs about 2 grown gates per such cell, where the pure growth model needs about 22.
\item \textbf{Bounds} (\cref{sec:bounds}, \cref{app:discrete,app:networks}). The ``short training $\Rightarrow$ small'' implication is a theorem for:
  \begin{itemize}
  \item the discrete model, for size read as depth, incoherence-weighted size, or readout norm (at depth one), though not for the raw gate count;
  \item diagonal linear networks;
  \item quadratic networks with fixed output signs;
  \item two-layer ReLU networks through the first escape.
  \end{itemize}
  It is a conjecture beyond the first escape. On the real networks, the only tight bounds are a-posteriori distillation certificates.
\item \textbf{Composition} (\cref{sec:composition}). Leap complexity depends only on Fourier supports. In a test whose predictions were recorded before any network result was inspected, intermediate features on a degree-4 term's own inputs make that term cheaper. They do so by $1.2$–$1.6\times$ in samples, and by up to $10\times$ in time, between targets of the same leap, and the gap grows with $d$. Features on other inputs do not help.
\end{enumerate}

\paragraph{Scope.} Everything empirical here is one architecture family, a $\mu$P two-hidden-layer ReLU MLP, mostly at width 128 and with Adam. Inputs are uniform on the hypercube, every target has leap at most 4 (parities of at most 4 bits), and the loss is square loss. We label every claim \status{Proved}, \status{Conditional}, \status{Conjecture}, \status{Heuristic} or \status{Numerical}.
